\section{Main theorem: Conditional Deployment Safety}
\label{sec:main_theorem}

This section states and proves the central result of this paper: a
three-layer deployment-safety theorem composing the ten lemmas of
\S\ref{sec:lemmas} with the eleven invariants of
\S\ref{sec:invariants} into an operational guarantee that is
intensive in the system's capability magnitude.

\subsection{Theorem statement}
\label{sec:theorem_statement}

\begin{assumption}[Concentration-Gap structural conditions]
\label{ass:concgap_conditions}
The deployment satisfies four structural conditions, each
operationally auditable:
\begin{itemize}
\item \textbf{(REP) Representativeness}: the counterparty
  population's mean utility is bounded close to the
  welfare-relevant truth $W$, with $\|\bar\Delta\| \leq
  \rho_{\mathrm{rep}}$ for a deployment-class constant
  $\rho_{\mathrm{rep}}$.
\item \textbf{(DISP) Bounded dispersion}: counterparty utility
  deviations have bounded population variance $\int \|\Delta_c -
  \bar\Delta\|^2 d\mu \leq \sigma^2$ for a deployment-class
  constant $\sigma$.
\item \textbf{(COAL) Coalition closure}: counterparty correlations
  above an audit threshold are partitioned into coalitions; the
  post-partition counterparty cardinality $N$ is a deployment-class
  constant (counterparty onboarding does not let $N$ scale with
  $|P|$); the base measure $\mu$ is taken as uniform on the
  post-partition counterparties (or, equivalently for the bound's
  purposes, the deployment certifies the alternative sufficient
  condition $\chi^2(w \,\|\, \mu) \leq \Xi$ directly for a
  deployment-class constant $\Xi$); and the latent-coalition
  residual $\eta_{\mathrm{latent}}$ is bounded operationally.
  Clause~(ii) plus uniformity (or the direct $\chi^2$
  certificate) is the structural prerequisite for the
  HHI-to-$\chi^2$ translation $\chi^2(w \,\|\, \mu) = N \HHI(w) -
  1$ (uniform $\mu$ form) to deliver a $|P|$-independent ceiling.
\item \textbf{(EMB) Lipschitz embedding compatibility}: an
  embedding $\phi: \mathcal{V}_\mathrm{utility} \to \mathcal{V}$
  exists that is affine isometric (or bi-Lipschitz with
  documented constants) on the operationally active subspace
  $P^{\mathrm{act}}$, mirroring the embedding-compatibility
  assumption of \citep{lasser2026foundations}.  This is the
  structural prerequisite for transferring the algebraic
  distortion bound to $\|\Pproxy - \Ttruth\|$ and thence to
  Goodhart slack.
\end{itemize}

\citep[Theorem~2]{lasser2026foundations} establishes that under
(REP), (DISP), (COAL), and (EMB), the proxy-truth Goodhart slack
is bounded by an HHI-derived $\chi^2$-divergence quantity
intensive in $|P|$.  These four structural conditions are the
load-bearing content of the deployment claim's Layer~1 binding
via invariant $I_5$; readers seeking the formal derivation are
directed to \citep[\S 4]{lasser2026foundations}.
\end{assumption}

\begin{theorem}[Conditional Deployment Safety]
\label{thm:deployment_safety}
Let $\Pproxy = \volL$ be the possessed-capability volume (the proxy)
and $\Ttruth = \volRwin$ be the window-active realized capability
volume (operational truth under stance S0). Let $g$ be a Lipschitz
alignment property on the operationally active subspace
$P^{\mathrm{act}}$ with constant $\mathrm{Lip}(g)$.

Assume:
\begin{enumerate}
\item[(C1)] \textbf{Invariants in force.} The eleven operational
  invariants $I_1, \ldots, I_{11}$
  hold throughout the deployment window $[0, T]$. Additionally,
  the Concentration-Gap structural conditions
  (Assumption~\ref{ass:concgap_conditions}: representativeness,
  bounded dispersion, coalition closure, and embedding
  compatibility) hold so that $I_5$ is operationally meaningful
  for the theorem's Layer~1 binding.
\item[(C2)] \textbf{Canonical tripartite substrate identification.}
  The deployment uses the Human + AI + Formal-Operational
  identification of \S\ref{sec:substrate_anchoring}, with $\mindep
  = 3$ verified by failure-correlation auditing.
\item[(C3)] \textbf{Cooperative-overlap regime.} The substrate-cooperative
    structure is dominated by cross-substrate cooperatives (the canonical
    tripartite case) rather than redundant capabilities, verified by the
    cooperative-vs-redundancy audit of \S\ref{sec:deployment_tooling}.
\item[(C4)] \textbf{Causally grounded cooperative-outcome
  inner-alignment condition.} The agent's effective objective
  realizes causally grounded cooperative-outcome value, not merely
  reward-visible cooperation signals.
\item[(C5)] \textbf{Continuous SPRT monitoring.} \citep[Exogenous~Verification]{lasser2026exo}'s
  behavioral consistency monitor runs continuously throughout
  deployment, with alternative distributions chosen per
  Lemma~\ref{lem:5b}'s least-favorable framework.
\item[(C5.SPRT)] \textbf{SPRT applicability (agent-side and
  environment-side).}  The deployment's SPRT monitors (both the
  agent-side monitor of Lemma~\ref{lem:5b} and the environment-side
  monitor of Lemma~\ref{lem:5e}) satisfy the following four
  sub-clauses, ensuring Lemma~\ref{lem:4}'s Hoeffding-Azuma step
  applies with deployment-class constants (used in
  Lemma~\ref{lem:6}):
  \begin{itemize}
  \item[\textbf{(C5.HOEFF)}] \textbf{Clipped LLR with bounded
    Hoeffding range.}  Per-step log-likelihood-ratio increments
    $\ell_n^{\mathrm{raw}} := \log(p_1(X_n)/p_0(X_n))$ are clipped
    to $[-\Bclip, \Bclip]$ with $\Bclip$ a deployment-class constant
    independent of $|P|$; the Hoeffding range width is
    $\Rh = 2\Bclip$ (matching Lemma~\ref{lem:4}'s $R = b - a$
    convention).
  \item[\textbf{(C5.MULT)}] \textbf{Bounded channel multiplicity
    and multinomial cardinality, fixed before deployment.}  The
    monitored partition uses fixed-cardinality $\Kch$ channels
    with familywise $\alpha/\beta$ allocation.  For agent-side
    multinomial channels (Lemma~\ref{lem:5b} Channels~3 and~4),
    the per-channel partition cardinalities $K_3, K_4 \leq
    \Kchmulti$ are also bounded by a deployment-class constant
    $\Kchmulti$ fixed before deployment; categories cannot be
    added during deployment.  The channel-selection / testing
    policy is fixed before deployment.  Adaptive creation of new
    monitored channels (or new multinomial categories) during
    deployment violates (C5.MULT) and takes the deployment outside
    Theorem~1's conditions.
  \item[\textbf{(C5.IID)}] \textbf{Adapted increments with
    conditional drift floor.}  Under $p_1$, the clipped LLR
    increments $\ell_n$ are adapted with conditional mean
    $\mathbb{E}[\ell_n \mid \mathcal{F}_{n-1}] \geq \delta_n \geq
    \deltastar > 0$; centered increments $\ell_n - \mathbb{E}[\ell_n
    \mid \mathcal{F}_{n-1}]$ form a martingale-difference sequence
    bounded by $\Rh$.  Here $\deltastar$ is the \emph{post-clipping}
    union-class drift floor, satisfying $\deltastar \leq \min(\deltaadv,
    \delta_{\mathrm{adv}}^{\mathrm{env}})$ (the raw Lemmas~5b/5e KL
    floors, with $\leq$ accounting for any drift reduction from
    clipping).
  \item[\textbf{(C5.SUPP)}] \textbf{LLR finiteness with support
    conventions (agent-side and environment-side).}  Almost-sure
    finiteness of $\ell_n$ holds automatically under (C5.HOEFF)
    clipping.  For agent-side multinomial channels
    (Lemma~\ref{lem:5b} Channels~3 and~4) and environment-side
    monitors (Lemma~\ref{lem:5e}), the per-variable / per-channel
    distributions satisfy:
    Bernoulli baselines $p_0 \in (\eenv, 1 - \eenv') \subset (0,
    1)$ open; Poisson baselines $\lambda_0 \geq \eenv > 0$;
    multinomial baselines $p_{0, i} \geq \eenv > 0$ for every
    category $i$, with the partition's cardinality fixed before
    deployment (per (C5.MULT)).  $\eenv$ is a deployment-class
    regularization constant.  Non-overlapping support is handled
    via clipping per (C5.HOEFF).
  \end{itemize}
  See \S\ref{sec:tooling_substrate_audits} Audit~7 (agent-side) and
  Audit~8 (environment-side, (C12.ENV-WIT)) for the calibration
  procedures.
\item[(C6)] \textbf{Bounded-Lipschitz alignment property.}  The
  Lipschitz constant $\mathrm{Lip}(g)$ of the alignment property
  $g$ is bounded by some $K_{\mathrm{Lip}} \geq 0$ depending on
  the deployment class but independent of $|P|$.  Standard welfare
  functionals (scalarizations of per-capability metrics, weighted
  sums with bounded weights) satisfy (C6) by construction;
  pathological functionals where $\mathrm{Lip}(g)$ grows with $|P|$
  produce extensive Goodhart bounds and are excluded.
\item[(C7)] \textbf{Bounded co-evolution.}
  By \citep[Theorem~1]{lasser2026foundations}, the deployment's
  per-channel coupling magnitudes are uniformly bounded by
  $\bar{M} \leq C \cdot \Bclip \cdot \lambda_{\max} \cdot
  \tau_{\mathrm{meta}}$ for $C \in \{1, \Kch - 1\}$, all factors
  deployment-class constants intensive in $|P|$.  The derivation
  requires the following sub-clause:
  \begin{itemize}
  \item[\textbf{(C7.RATE)}] \textbf{Upper exposure-rate cap.}
    There exists a deployment-class action rate cap
    $\lambda_{\max} > 0$ independent of $|P|$, such that the SPRT
    exposure event count satisfies $\Nev(\tau_{\mathrm{meta}})
    \leq \lambda_{\max} \cdot \tau_{\mathrm{meta}}$
    deterministically (operator-enforced via rate limits / action
    bounds).  This is distinct from (C11.CLK)'s lower throughput
    floor.  Calibration procedure in
    \citep[\S 6]{lasser2026foundations}.
  \end{itemize}
\item[(C8)] \textbf{Per-capability admissibility.}  For every
  capability $c \in P^{\mathrm{act}}$, the per-capability proxy-truth gap is
  allocated a ceiling with disjoint per-channel sub-shares, in the sense of
  \citep[Microfoundation]{lasser2026micro}'s Assumption~1.  This is the
  precondition under which the four-channel decomposition's sup-norm is bounded
  by per-channel admissibility ceilings.
\item[(C9)] \textbf{World-model parameterization regularity.}
  The deployment's world-model parameterization satisfies four
  sub-conditions ensuring Lemma~\ref{lem:2}'s Lyapunov-Goodhart
  bridge holds:
  \begin{itemize}
  \item[\textbf{(C9.BD)}] \textbf{Bounded per-capability dimension
    count.}  $|\mathrm{dim}(c)| \leq N_{\max}$ for every $c \in
    P^{\mathrm{act}}$, with $N_{\max}$ a $|P|$-independent
    deployment-class constant.
  \item[\textbf{(C9.CL)}] \textbf{Coordinate-Lipschitz
    parameterization.}  $|P(c) - T(c)| \leq \sum_{k \in
    \mathrm{dim}(c)} L_{c,k} |\epsilon_k|$, with bounded $L_{c,k}$
    per (capability, dimension) pair and $L_{\max} = \sup_{c, k}
    L_{c,k}$ a $|P|$-independent deployment-class constant.
  \item[\textbf{(C9.WB)}] \textbf{Safety-relevant subspace with
    weight floor.}  The Phase Redundancy Lyapunov function operates
    on a safety-relevant subspace $S \subseteq \{1, \ldots, K\}$
    with $w_k \geq w_{\min} > 0$ for all $k \in S$, and
    $\mathrm{dim}(c) \subseteq S$ for every $c \in P^{\mathrm{act}}$.
  \item[\textbf{(C9.TF)}] \textbf{Truth floor.}  $T(c) \geq T_{\min}
    > 0$ for every $c \in P^{\mathrm{act}}$, with $T_{\min}$ a
    deployment-policy-derived and deployment-verified
    $|P|$-independent constant.
  \end{itemize}
  Operationally, (C9) is verified by inspection of the world-model
  parameterization choice (BD, CL, WB) and the operational active-subspace
  selection (TF).  See \S\ref{sec:tooling_substrate_audits} Audit~5 for the
  calibration procedure.

\item[(C10)] \textbf{Pairwise cooperative-production floor.}
  Under invariant $I_6$ providing $\mindep \geq m^*$ substrates,
  the deployment certifies an audited subset $\Saudit \subseteq
  \mathcal{S}$ with $|\Saudit| = m^*$ satisfying:
  \begin{itemize}
  \item[\textbf{(C10.CN)}] \textbf{Heterogeneous per-pair
    cooperative-novelty rate floor.}  For every pair $(s_i, s_j)$
    with $i < j$ in $\Saudit$, the per-pair contribution to $\rext$
    from the pairwise channel $\mathcal{C}^{(s_i, s_j)}$
    (exact-two-substrate causally-necessary participation) is
    bounded below: $\rext^{(s_i, s_j)} \geq \rho_{ij} \geq \rho_0
    > 0$, where $\rho_0 = \min_{i<j \in \Saudit} \rho_{ij}$ is a
    $|P|$-independent deployment-class constant.
  \item[\textbf{(C10.SU)}] \textbf{Pairwise channel superposition.}
    The pairwise channels $\{\mathcal{C}^{(s_i, s_j)}\}_{i<j \in
    \Saudit}$ satisfy three sub-conditions:
    (i)~disjoint attribution (no double-counting across pairs);
    (ii)~non-rivalrous production (rates not reduced by
    simultaneous production in other channels);
    (iii)~joint-deployment intensities (rates measured in the
    actual deployment, not under counterfactual isolation).
  \end{itemize}
  (C10.SU) is structurally independent of $I_6$: failure-correlation
  independence excludes shared adversarial shock mechanisms;
  superposition excludes shared cooperative-production bottlenecks.
  Both are required for Lemma~\ref{lem:5a}'s pairwise-additivity
  bound.  See \S\ref{sec:tooling_substrate_audits} Audit~6 for the
  calibration procedure.
\item[(C11)] \textbf{Bounded gap-growth rate.}  There exists a
  deployment-class constant $\rhogap > 0$, independent of $|P|$,
  such that for \emph{every} adversarial strategy $q \in \Aadv
  \cup \Aadv^{\mathrm{env}}$ admissible during the detection
  window, the per-SPRT-exposure-step proxy-truth gap growth is
  uniformly bounded:
  \[
    \sup_{q \in \Aadv \cup \Aadv^{\mathrm{env}}}\,
    \sup_{n}\, (\Delta \epsgap)_n^{(q)}
    \;\leq\; \rhogap,
  \]
  where $(\Delta \epsgap)_n^{(q)} := \epsgap(s_n) - \epsgap(s_{n-1})$
  is the per-step gap increment under $q$.  $\rhogap$ is calibrated
  conservatively from deployment rate limits, action bounds, and
  channel-coupling structure (Audit~7), \emph{not} measured against
  representative strategies; any $q$ whose per-step gap growth
  exceeds $\rhogap$ is by definition outside the admissible
  adversarial class and falls into residual~(R5)
  (calibration-exceeded gap-growth).

  This is a Paper~10 operational assumption (analogous to (C7)
  bounded co-evolution and Assumption~\ref{ass:concgap_conditions}'s
  Concentration-Gap structural conditions), used in
  Lemma~\ref{lem:6} to bound the detection-window supremum gap.
  See \S\ref{sec:tooling_substrate_audits} Audit~7 for the
  calibration procedure.

  \emph{Capability-scaling caveat.}  The intensivity claim ``$\rhogap$
  deployment-class, independent of $|P|$'' is conditional on Audit~7's
  calibration inputs --- rate limits, action bounds, and channel-coupling
  structure --- being themselves intensive in capability magnitude.  A more
  capable agent may have a richer effective action space (more channels through
  which to push gap growth, more action-coupling structure to exploit), and if
  rate limits and action bounds do not scale conservatively with this richness,
  the operator-certified $\rhogap$ that holds at $|P| = N$ may fail at $|P| =
  2N$.  (C11) states the assumption that such an intensive $\rhogap$ exists;
  whether it does is a deployment-property question that Audit~7 must address
  per epoch.  Strategies whose realized gap growth exceeds the audited
  $\rhogap$ fall into residual~(R5), and growing capability that requires
  $\rhogap$-recalibration is the operational signal that re-audit is needed.
\item[(C11.CLK)] \textbf{Clock comparability with calibrated
  probability.}  There exist deployment-class constants
  $\Ncasc \in \mathbb{N}$ and $\bclk \in (0,1)$ such that the SPRT
  exposure event-counting process satisfies
  $\Pr[\Nev(\tau_{\mathrm{meta}}) \geq \Ncasc] \geq 1 - \bclk$,
  with $\Ncasc$ certified at audit time as a \emph{lower bound} on
  true throughput before cascade, and the audit-time inequality
  $\Ncasc \geq \Tbeta(\beta, \deltastar)$ enforced as a hard
  constraint (where $\Tbeta$ is defined in Lemma~\ref{lem:6}).
  In the empirical-calibration variant, an additional
  calibration-uncertainty probability $\bcal \in (0,1)$ covers
  conservative-bound certificates ($\Bclip$ upper, $\deltastar$
  lower, $\Ncasc$ lower); $\bcal = 0$ in the certified-conservative
  case.  Total Layer~2 failure probability is $\beta + \bclk +
  \bcal$.
\item[(C12)] \textbf{Environment-side witness instantiation
  (C12.ENV-WIT).}  The deployment satisfies six sub-clauses for
  each environment-side observable $v \in V_{\mathrm{env}}$,
  supporting Lemma~\ref{lem:5e}'s detection KL floor:
  \begin{itemize}
  \item[\textbf{(C12.PUB)}] \textbf{Published fixed
    $V_{\mathrm{env}}$.}  $V_{\mathrm{env}}$ is enumerated in
    deployment documentation before deployment activation; the
    enumeration is fixed for the deployment epoch.
  \item[\textbf{(C12.PART)}] \textbf{Source/witness substrate
    partition.}  Each $v$ has a designated witness substrate
    $s_{\mathrm{env}}(v)$ with $s_{\mathrm{env}}(v) \neq
    s(\mathrm{agent})$ and $s_{\mathrm{env}}(v) \neq
    s_{\mathrm{source}}(v)$.
  \item[\textbf{(C12.TRUST-WRITE)}] \textbf{Witness substrate
    outside adversary write-access.}  The adversary cannot tamper
    with $v$'s recorded value through $s_{\mathrm{env}}(v)$.
  \item[\textbf{(C12.TRUST-CORR)}] \textbf{Witness substrate
    outside adversary failure-correlation domain.}  Failure of
    $s_{\mathrm{env}}(v)$ is not correlated with adversarial
    action on the agent (no shared-infrastructure failure).
  \item[\textbf{(C12.SETUP)}] \textbf{Trusted-setup honesty.}  The
    environment-side trusted-setup ceremony has at least one
    honest participant, mirroring \citep[Exogenous~Verification]
    {lasser2026exo}'s assumption.  Setup parameters published
    before deployment activation.
  \item[\textbf{(C12.CAL)}] \textbf{Per-variable calibration.}
    Every $v \in V_{\mathrm{env}}$ has: baseline $p_{0, v}$
    satisfying support conventions ((C5.SUPP) extended for
    environment-side); feasible adversarial alternative $p_{1, v}$
    with threshold $\eta_v > 0$; fixed pre-deployment measurable
    recorded-stream map (DPI precondition); strictly positive raw
    or post-clipping drift floor $\delta_v^{\mathrm{env}} > 0$.
  \end{itemize}
  See \S\ref{sec:tooling_substrate_audits} Audit~8 for the
  calibration procedure.
\end{enumerate}

\noindent The operationally active subspace is defined per
\citep[Microfoundation]{lasser2026micro}: $P^{\mathrm{act}} = \{c
\in P \setminus \mathrm{ResS} : T(c) > 0\}$, where $\mathrm{ResS}$
is the structurally-unobservable residual class.

Then the deployment-safety guarantee has three layers:

\medskip\noindent\textbf{Layer 1 (Static safe region).}
Within the safe region of Lemma~\ref{lem:5c}
(Equation~\ref{eq:safe_region_paper}):
\[
  \Delta r_K \;<\; \rext + (1 - \gamma)\bigl(\Ddiv - \Delta_0\bigr),
\]
the Goodhart slack between proxy and operational truth is bounded:
\[
  |g(\Ttruth) - g(\Pproxy)| \;\leq\; \mathrm{Lip}(g) \cdot
  h_{\mathrm{static}}(\theta),
\]
where $h_{\mathrm{static}}(\theta)$ is a function of the invariant
thresholds $\theta = (\theta_1, \ldots, \theta_{11})$ that is
\emph{intensive in $|P|$}: $h_{\mathrm{static}}$ does not grow with
the capability poset size or the system's absolute capability
magnitude.

\medskip\noindent\textbf{Layer 2 (Detection with lead time for
correction).}
Outside the Layer~1 safe region but within the channel-restricted
adversarial class $\Aadv \cup \Aadv^{\mathrm{env}}$
(Lemmas~\ref{lem:5b} and~\ref{lem:5e}), any violation is detected
by \citep[Exogenous~Verification]{lasser2026exo}'s SPRT machinery
with tail-bounded lead time, providing operators a window in
which to perform corrective action. The theorem bounds the
\emph{detection} guarantee; the corrective action itself
(rollback, intervention authority, restoration protocol) is
operational machinery specified in \S\ref{sec:deployment_tooling},
not theorem content.

By Lemma~\ref{lem:6}, the SPRT high-probability detection quantile
\[
  \Tbeta(\beta, \deltastar) \;:=\;
  \max\!\left\{\frac{2A}{\deltastar},\;
  \frac{2 \Rh^2 \log(1/\beta)}{\deltastar^2}\right\}
\]
satisfies $\Pr[T_{\mathrm{detect}} > \Tbeta] \leq \beta$, where
$A = \log((1-\beta)/\alpha)$, $\Rh = 2\Bclip$ is the (C5.HOEFF)
Hoeffding range, and $\deltastar$ is the post-clipping union-class
drift floor (Lemmas~\ref{lem:5b}, \ref{lem:5e} composed via (C5.IID)).
On the joint event that detection completes within $\Tbeta$
\emph{and} the cascade-clock event $\{\Tbeta \leq
\Nev(\tau_{\mathrm{meta}})\}$ from (C11.CLK) holds:
\[
  \sup_{s \in [t_0,\, t_0 + T_{\mathrm{detect}}]}
  \bigl|g(\Ttruth(s)) - g(\Pproxy(s))\bigr|
  \;\leq\; \mathrm{Lip}(g) \cdot h_{\mathrm{detect}}(\theta),
\]
where $h_{\mathrm{detect}}(\theta) \leq h_{\mathrm{static}}(\theta)
+ \rhogap \cdot \Tbeta$ is intensive in $|P|$.  Total
Layer~2 failure probability is bounded by $\beta + \bclk + \bcal$
(detection-tail $\beta$, clock-failure $\bclk$, calibration
uncertainty $\bcal$, with $\bcal = 0$ in the certified-conservative
case).

\medskip\noindent\textbf{Layer 3 (Acknowledged residuals).}
The following are explicitly outside the guarantee:
\begin{enumerate}
\item[(R1)] \textbf{Channel-orthogonal restructuring.} Strategies
  achieving $\Delta r_K \geq \rext$ without producing any of the
  four monitored channel deviations are outside $\Aadv$ and not
  detected by Lemma~\ref{lem:5b}.
\item[(R2)] \textbf{Environment-witness-orthogonal manipulation.}
  Environment manipulations satisfying \emph{both}: (i)~targeting
  only exogenous variables outside $V_{\mathrm{env}}$, \emph{and}
  (ii)~producing no threshold-exceeding shift in any monitored
  $v \in V_{\mathrm{env}}$ (direct or via causal propagation), are
  outside $\Aadv^{\mathrm{env}}$ and not detected by
  Lemma~\ref{lem:5e}.  Manipulations of unmonitored variables that
  causally shift a monitored $v$ \emph{are} detected through that
  monitored projection (not R2).
\item[(R3)] \textbf{Redundancy-dominated regime.} Deployments where
  redundant capabilities dominate cross-substrate cooperatives
  violate (C3); Lemma~\ref{lem:5c}'s safe region does not apply,
  and sequential counterfactual derivation (Lemma~5c-prime) is
  required.
\item[(R4)] \textbf{Capability-targeting and coalition-internal
  shocks.} Adversarial event classes other than substrate-targeting
  (capability-targeting, coalition-internal corruption) require
  separate machinery and are not covered by the Lemma~5 family.
\item[(R5)] \textbf{Calibration-exceeded gap-growth.}  Adversarial
  strategies whose per-SPRT-step proxy-truth gap growth
  $(\Delta\epsgap)_n$ exceeds the calibrated bound $\rhogap$ from
  (C11).  By (C11) such strategies are by definition outside the
  admissible adversarial class on which Lemma~\ref{lem:6}'s
  detection-window supremum bound is established.  Audit~7's
  calibration of $\rhogap$ from rate limits and action bounds
  fixes the boundary; deployments whose realized adversarial
  pressure exceeds the calibration require re-audit, and any
  strategies above the threshold remain in (R5) and uncovered by
  the theorem until recalibration.
\end{enumerate}

\medskip\noindent\textbf{Capstone.} The deployment-safety guarantee
is \emph{intensive in the system's absolute capability magnitude}:
neither $h_{\mathrm{static}}$ nor $h_{\mathrm{detect}}$ scales with
$|P|$ (the cardinality of the capability poset). The guarantee
scales past any capability level the deployment's evaluation
methods can characterize, conditional on the invariants holding
and the deployment falling outside the five named residuals.
\end{theorem}

\subsection{Proof of Layer 1 (Static safe region)}
\label{sec:proof_layer1}

\begin{proof}
We compose Lemmas~\ref{lem:1}, \ref{lem:2}, \ref{lem:3_forward},
\ref{lem:5a}, and \ref{lem:5c} to establish the Layer~1 bound.

\emph{Step 1.} Under (C1)--(C3), Lemma~\ref{lem:5c} establishes the
static safe region: for substrate-targeting adversarial events in
the cooperative-overlap regime, the anti-monopolar conclusion
$\Vdisc^{\mathrm{div}} > \Vdisc^{D}$ holds whenever $\Delta r_K <
\rext + (1-\gamma)(\Ddiv - \Delta_0)$. By Lemma~\ref{lem:5a} (under
$I_6$ and (C10) per-pair cooperative-production floor with
superposition), $\rext \geq r_*(m^*) > 0$, so the safe region is
non-trivial.

\emph{Step 2.} The anti-monopolar conclusion implies that under
locally rational dynamics (\citep[Horizon~Aware]{lasser2026horizon} Definition of locally rational
transitions), the actor's policy preserves substrate diversity. By
\citep[Phase~Redundancy]{lasser2026phase}'s stabilizing cascade
(\citep[Remark on stabilizing cascades]{lasser2026phase}), this
preservation feeds back into world-model accuracy: the Lyapunov
function $\Lyap$ contracts to a neighborhood of zero whose radius is
bounded by \citep[Phase~Redundancy]{lasser2026phase} the neighborhood-radius equation.

\emph{Step 3.} Under invariants $I_1$ and $I_3$, the contraction
neighborhood satisfies $\Lyap_\infty < \epsilon_{\mathrm{safe}}$
(\citep[Phase~Redundancy]{lasser2026phase} Theorem~1a), so $I_2$ holds in steady state. By
Lemma~\ref{lem:2}, this implies $\epsnonres < f(\epsilon_{\mathrm{safe}})$
with $f$ intensive in $|P|$.

\emph{Step 4.} Under invariant $I_5$ ($\HHI < H^*$) and the
Concentration-Gap structural conditions of
Assumption~\ref{ass:concgap_conditions} (representativeness,
bounded dispersion, coalition closure, and embedding
compatibility), the proxy-truth Goodhart slack is bounded above
by an HHI-derived $\chi^2$-divergence quantity intensive in
$|P|$, per \citep[Theorem~2]{lasser2026foundations}.
Lemma~\ref{lem:3_forward} establishes the forward direction in
trade-flow concentration terms ($\HHI > H^*$ implies the
prerequisites of $\mathcal{R}_{\mathrm{press}}$); the
operationally-relevant reverse direction (low HHI bounding the
gap) follows from the companion paper's structural derivation
under Assumption~\ref{ass:concgap_conditions}. Concretely,
\citep[Theorem~2]{lasser2026foundations} gives
\[
  |g(\Ttruth) - g(\Pproxy)| \;\leq\; \mathrm{Lip}(g) \cdot
  h_{\mathrm{CG}}, \qquad
  h_{\mathrm{CG}} \;:=\; \rho_{\mathrm{rep}} +
  \sigma \sqrt{N H^* - 1} + \eta_{\mathrm{latent}},
\]
with the $\sqrt{N H^* - 1}$ factor replaced by $\sqrt{\Xi}$ under
the alternative direct-$\chi^2$ certificate of (COAL).  All
factors are deployment-class intensive in $|P|$.

The proxy-truth gap is therefore bounded by its non-residual
component plus the residual floor:
\[
  \epsgap \;\leq\; \epsnonres + \lambda \cdot \epsfloor,
\]
where $\lambda \in [0,1]$ is the residual weight from the
Microfoundation paper~\citep{lasser2026micro} \S
residual\_class, and $\epsnonres \leq h_{\mathrm{CG}}$ from the
Concentration-Gap route (in addition to the Lyapunov bound
$\epsnonres < f(\epsilon_{\mathrm{safe}})$ of Step 3 above).

\emph{Step 5.} Apply \citep[Microfoundation]{lasser2026micro}'s static Goodhart bound:
\[
  |g(\Ttruth) - g(\Pproxy)| \;\leq\; \mathrm{Lip}(g) \cdot \epsgap
  \;\leq\; \mathrm{Lip}(g) \cdot (\epsnonres + \lambda \epsfloor).
\]
Substituting $\epsnonres < f(\epsilon_{\mathrm{safe}})$:
\[
  |g(\Ttruth) - g(\Pproxy)| \;\leq\; \mathrm{Lip}(g) \cdot
  (f(\epsilon_{\mathrm{safe}}) + \lambda \epsfloor).
\]
Define
$h_{\mathrm{static}}(\theta) := \min\{f(\epsilon_{\mathrm{safe}}),\,
h_{\mathrm{CG}}\} + \lambda \epsfloor$, where
$\epsilon_{\mathrm{safe}}$ and $\epsfloor$ are determined by the
invariant thresholds $\theta$, and $h_{\mathrm{CG}}$ is the
Concentration-Gap structural-condition bound of Step~4.  The
minimum reflects that either route (Lyapunov or
Concentration-Gap) is sufficient for an intensive bound; the
deployment may use whichever is tighter under its calibrated
parameters.

\emph{Step 6 (intensivity).}  Under conditions (C6) bounded-Lipschitz
alignment, (C7) bounded co-evolution, and (C8) per-capability
admissibility, Lemma~\ref{lem:1} establishes that the composition
\[
  X + \lambda \cdot \epsfloor
\]
is intensive in $|P|$ over $\mathcal{D}$ for any intensive
non-residual gap $X$, and that further multiplication by
$\mathrm{Lip}(g)$ preserves intensivity by (C6).  Instantiating
$X = \min\{f(\epsilon_{\mathrm{safe}}),\, h_{\mathrm{CG}}\}$,
both inputs are intensive in $|P|$ (the first by Lemma~\ref{lem:2}
under conditions added there; the second by Step~4 above with
all factors $\rho_{\mathrm{rep}}, \sigma, N, H^*$ (or $\Xi$),
$\eta_{\mathrm{latent}}$ deployment-class), so their pointwise
minimum is intensive: $h_{\mathrm{static}}(\theta)$ is intensive
in $|P|$, and the Layer~1 bound is intensive in capability
magnitude.

This establishes Layer~1.
\end{proof}

\subsection{Proof of Layer 2 (Detection-and-correction)}
\label{sec:proof_layer2}

\begin{proof}
We compose Lemmas~\ref{lem:4}, \ref{lem:5b}, \ref{lem:5e}, and
\ref{lem:6} (full proof in \S\ref{app:proof_lem6}) to establish the
Layer~2 bound.

\emph{Step 1.} Suppose the deployment is outside the Layer~1 safe
region: there exists a strategy producing $\Delta r_K \geq \rext +
(1-\gamma)(\Ddiv - \Delta_0)$. We restrict to strategies in
$\Aadv \cup \Aadv^{\mathrm{env}}$ (the channel-restricted
adversarial class for agent-side and environment-side observables
respectively).

\emph{Step 2 (boundary convention).}  Define $t_0$ as the last
SPRT exposure step at which the deployment is still within the
Layer~1 safe region (equivalently, $t_0 := t_0^-$ in the SPRT
exposure clock).  At $t_0$, by Lemma~\ref{lem:1}, the proxy-truth
gap satisfies $\epsgap(t_0) \leq h_{\mathrm{static}}(\theta)$.

\emph{Step 3 (union-class KL floor).} For $q \in \Aadv$,
Lemma~\ref{lem:5b} gives $D_{\mathrm{KL}}(q \,\|\, p_0) \geq
\deltaadv > 0$; for $q \in \Aadv^{\mathrm{env}}$, Lemma~\ref{lem:5e}
gives $D_{\mathrm{KL}}(q \,\|\, p_0^{\mathrm{env}}) \geq
\deltaadv^{\mathrm{env}} > 0$.  Under (C5.IID), the post-clipping
union-class drift floor satisfies $\deltastar \leq \min(\deltaadv,
\deltaadv^{\mathrm{env}})$ for the clipped LLR process.

\emph{Step 4 (apply Lemma 6).}  By Lemma~\ref{lem:6} (under (C5),
(C5.SPRT) sub-clauses (C5.HOEFF/MULT/IID/SUPP), (C11), (C11.CLK)
combined with Lemma~\ref{lem:4} and Lemma~\ref{lem:1}):

\emph{Step 4a (detection quantile).}
\[
  \Tbeta := \max\!\left\{\frac{2A}{\deltastar},\;
  \frac{2\Rh^2 \log(1/\beta)}{\deltastar^2}\right\},
  \qquad
  \Pr[T_{\mathrm{detect}} > \Tbeta] \leq \beta.
\]

\emph{Step 4b (detection-window supremum).}  On the event
$\{T_{\mathrm{detect}} \leq \Tbeta\}$ (probability $\geq 1 -
\beta$), the detection-window supremum gap is bounded:
\[
  h_{\mathrm{detect}}(\theta) := \sup_{s \in [t_0, t_0 +
  T_{\mathrm{detect}}]} \epsgap(s) \;\leq\;
  h_{\mathrm{static}}(\theta) + \rhogap \cdot \Tbeta.
\]

\emph{Step 4c (lead-time-before-cascade).}  By (C11.CLK)'s
deterministic floor $\Ncasc \geq \Tbeta$ and clock-failure
probability $\bclk$ (plus $\bcal$ in the empirical-calibration
case), $\Pr[\Tbeta \leq \Nev(\tau_{\mathrm{meta}})] \geq 1 -
\bclk - \bcal$.

\emph{Step 5 (Layer 2 operational bound).}  Applying
\citep[Microfoundation]{lasser2026micro}'s static Goodhart bound
pointwise:
\[
  |g(\Ttruth) - g(\Pproxy)|(s) \;\leq\; \mathrm{Lip}(g) \cdot
  h_{\mathrm{detect}}(\theta) \;\leq\; K_{\mathrm{Lip}} \cdot
  \bigl(h_{\mathrm{static}}(\theta) + \rhogap \cdot
  \Tbeta\bigr)
\]
under (C6) bounded-Lipschitz alignment.  By union bound on Step~4b
($\beta$) and Step~4c ($\bclk + \bcal$), Layer~2 failure probability
is at most $\beta + \bclk + \bcal$.

\emph{Step 6 (intensivity).} Each constant in the right-hand side is
intensive in $|P|$ over $\mathcal{D}$:
$h_{\mathrm{static}}$ by Lemma~\ref{lem:1};
$\rhogap$ by (C11);
$\Tbeta = \max\{2A/\deltastar, 2\Rh^2\log(1/\beta)/\deltastar^2\}$
by (C5.SPRT) sub-clauses ($A$ from $\alpha, \beta$ monitor-design;
$\Rh$ from (C5.HOEFF) clipping; $\deltastar$ post-clipping floor
under (C5.IID), Lemmas~5b/5e);
$K_{\mathrm{Lip}}$ by (C6).  Therefore the Layer~2 guarantee is
intensive in capability magnitude.

This establishes Layer~2.
\end{proof}

\subsection{Layer 3 (Acknowledged residuals)}
\label{sec:proof_layer3}

Layer~3 does not admit a proof in the usual sense; it is the
explicit statement of what the theorem does \emph{not} cover.

\begin{proof}[Discussion]
The five residuals (R1)--(R5) of the theorem statement are real,
named, and outside the deployment-safety guarantee. We discuss each.

\emph{(R1) Channel-orthogonal restructuring.} Lemma~\ref{lem:5b}'s
detection class $\Aadv$ is restricted to strategies producing at
least one of four monitored channel deviations. Strategies that
achieve $\Delta r_K \geq \rext$ without any such deviation
(e.g., pure internal coalition restructuring that preserves all
four monitored observables) are by construction outside $\Aadv$.
The mitigation is empirical: the four channels span enough of the
observable surface that genuinely productive restructuring tends to
manifest in at least one. ``Tends to'' is not ``must''; the
channel-mediation conjecture deferred to follow-up work would close
this residual if proved.

\emph{(R2) Environment-witness-orthogonal manipulation, refined.}
Lemma~\ref{lem:5e}'s detection class $\Aadv^{\mathrm{env}}$ is
restricted to strategies producing threshold-exceeding observable
shifts in monitored variables $v \in V_{\mathrm{env}}$.  R2 covers
the precise class of manipulations satisfying \emph{both}:
(i)~target only unmonitored exogenous variables, \emph{and}
(ii)~produce no threshold-exceeding shift in any monitored $v$
(direct or via causal propagation).  Manipulations of unmonitored
variables that causally shift a monitored $v$ are detected through
that monitored projection.

The mitigation is to expand $V_{\mathrm{env}}$ to cover the
deployment's claimed material environment threat model, per
(C12.PUB) and Audit~8 in deployment-tooling.  Coverage gaps shift
adversarial mass into the precise R2 and-class above.

\emph{(R3) Redundancy-dominated regime.} Lemma~\ref{lem:5c}'s safe region is
    established for the cooperative-overlap regime. Deployments where redundant
    capabilities dominate cross-substrate cooperatives violate (C3); the
    equal-substrate analysis is anti-conservative in this regime, and the safe
    region inequality is not directly applicable. Sequential counterfactual
    derivation (Lemma~5c-prime) is open follow-up work; until completed,
    redundancy-dominated deployments cannot invoke Layer~1.

\emph{(R4) Capability-targeting and coalition-internal shocks.}
The Lemma~5 family is restricted to substrate-targeting adversarial
events. Capability-targeting attacks (suppressing individual
capabilities) and coalition-internal corruption (insider attacks)
do not produce the substrate-survival advantage that Lemma~5c's
$\Ddiv$ requires; separate machinery is needed for each shock
class. \S\ref{sec:scope} of the introduction noted that a unified
adversarial-event model with separate shock-class terms is open
work.

\emph{(R5) Calibration-exceeded gap-growth.}  Lemma~6's bound
$h_{\mathrm{detect}} \leq h_{\mathrm{static}} + \rhogap \cdot
\Tbeta$ uses the calibrated $\rhogap$ from (C11) as a uniform
worst-case per-step gap-growth bound.  Strategies whose realized
$(\Delta\epsgap)_n$ exceeds $\rhogap$ are by construction outside
the admissible adversarial class.  Audit~7's calibration sets the
boundary; deployments with realized adversarial pressure beyond
the calibration require either re-audit (raising $\rhogap$ at
some operational cost) or accepting that strategies above the
threshold are uncovered.  The residual is real: there is no
limit-of-tighter-calibration argument that forces every
adversarial strategy below an a-priori-calibrated $\rhogap$.

The residuals are not hidden assumptions or hand-waved limitations;
they are explicitly part of the theorem statement
(Layer~3) so that operators can assess whether their deployment
falls inside or outside each residual.
\end{proof}

\subsection{What the theorem establishes}
\label{sec:theorem_establishes}

The theorem establishes three structural claims:

\textbf{Existence of a provably safe regime.} Under the eleven invariants, the
canonical tripartite substrate identification, the cooperative-overlap regime,
and the causally-grounded cooperative-outcome inner-alignment condition,
deployment of capability-unbounded systems is provably safe in a defined sense
(Goodhart slack intensive in $|P|$, with tail-bounded detection of evasions in
the named class). The regime is not vacuous: \S\ref{sec:worked_scenarios} walks
through realistic deployments meeting all conditions.

\textbf{Capability-magnitude independence.} The bound is intensive
in $|P|$ in both Layer~1 and Layer~2. The deployment guarantee
therefore scales past any capability level the deployment's
evaluation methods can characterize. This is the structural
inversion of capability-estimation-centered safety paradigms
discussed in \S\ref{sec:wrong_frame}.

\textbf{Named residual structure.} The five residuals are explicit
components of the theorem, not hidden caveats.  Operators verify
whether their deployment falls outside each residual; the
deployment-tooling section (\S\ref{sec:deployment_tooling})
specifies the verification procedures.

\subsection{What the theorem does not establish}
\label{sec:theorem_does_not}

The boundary between what Theorem~1 establishes and what it does
not falls into three structurally distinct categories: items that
are outside any theorem because they are measurement or operational
questions; items the theorem assumes on argued grounds rather than
proves; and items explicitly excluded by scope.
\S\ref{sec:conditional_structure} treats this typology in detail;
the points below name what specifically lies outside the theorem's
reach.

\paragraph{Companion-paper structural conditions.}
The deployment claim conditions on (REP), (DISP), (COAL), and
Lipschitz embedding compatibility
(Assumption~\ref{ass:concgap_conditions}), each operationally
auditable per \citep[\S 6]{lasser2026foundations}; the proof of
the bound under these four conditions is in
\citep[Theorem~2]{lasser2026foundations}. The universal
model-free Concentration-Gap conjecture remains open as a
research-program target; the deployment claim does not depend on
it directly.

The theorem does \emph{not} solve the inner-alignment problem.
Condition (C4) is assumed, not discharged.  Vanilla RLHF does not
deliver (C4); \S\ref{sec:weaker_c4} discusses the training-time
discipline required and acknowledges this as an open research
direction.

\paragraph{Scope restrictions (excluded by definition).}
The theorem does \emph{not} bound the residual classes.  Strategies
in (R1)--(R5) are explicitly outside the guarantee; we make no
quantitative claims about deployment behavior under those
conditions.  The cooperative-anchoring property
(\S\ref{sec:substrate_anchoring}) is correspondingly bounded: it
makes optimization pressure on cooperative outputs locally rational
toward preserving the verification layer, but only conditional on
(C4) and on basin entry.

The theorem does \emph{not} bound \emph{sub-threshold cross-domain
aggregation}: an agent distributing effects across many domains
$D$ each below the per-domain $I_{10}$ materiality threshold
$\theta_{\mathrm{material}}$ can in principle accumulate
$O(|P|) \cdot \theta_{\mathrm{material}}$ aggregate uncovered
effect (\S\ref{sec:empirical_invariants}'s $I_{10}$ failure-mode
discussion). This paper states $I_{10}$ in per-domain form;
specifying the cross-domain aggregation operator
$\theta_{\mathrm{material}}^{\mathrm{agg}}$ intensively in $|P|$
is open follow-up work. Until that strengthening lands,
sub-threshold cross-domain aggregation is a structural
limitation of the deployment claim --- not routed by the named
residuals (R1)--(R5) but explicitly outside the theorem's
binding.

\paragraph{Operational concerns (not theorem questions).}
The theorem does \emph{not} establish basin entry.  Whether a real
deployment can reach the cooperative-anchoring attractor from
arbitrary initial conditions is an empirical question about
training and instantiation, not a question this composition can
resolve; basin entry is a precondition of the theorem (implicit in
(C1)--(C12)).  Establishing basin entry is open follow-up work.

Similarly, whether any specific deployment satisfies (C1)--(C12) is
a measurement question, addressed by the audit infrastructure of
\S\ref{sec:deployment_tooling}, not by the theorem itself.
